Convex optimization is tractable because local derivative information controls global behavior.
The definitions and inequalities below formalize the properties we will use repeatedly; see \citet[Ch.~9]{boyd_vandenberghe_2004_convex}.

\begin{definition}[Convex function (first-order condition)]
Let $S \subseteq \R^d$ be convex, and let $f:S\to\R$ be differentiable.
We say $f$ is \emph{convex} if for all $x,y \in S$,
\begin{equation}
  f(y) \ge f(x) + \nabla f(x)^\top (y-x).
  \label{eq:convex-first-order}
\end{equation}
\end{definition}

Geometrically, \eqref{eq:convex-first-order} says the graph of $f$ lies above every tangent hyperplane.
Convexity does not require differentiability (one can state analogous conditions using subgradients), but in these notes we mostly assume differentiability to keep the calculus clean.

\begin{figure}[t]
  \centering
  \begin{tikzpicture}
  \begin{axis}[
    width=0.8\linewidth,
    height=5.0cm,
    xmin=-1.1, xmax=1.1,
    ymin=-0.1, ymax=0.9,
    axis lines=left,
    ticks=none,
    xlabel={$x$},
    ylabel={$f(x)$},
    clip=false,
  ]
    % A convex function (representative shape).
    \addplot[black,thick,domain=-1.1:1.1,samples=200] {0.15 + 0.4*x^2 + 0.1*x^4};

    % Tangent at x0=0.6.
    \addplot[stat221Red,thick,domain=-1.1:1.1,samples=2]
      {0.30696 + 0.5664*(x-0.6)};

    % Mark x0 and another point y.
    \addplot[only marks,mark=*,mark size=1.6pt] coordinates {(0.6,0.30696) (-0.4,0.21656)};
    \node[anchor=north] at (axis cs:0.6,0) {$x$};
    \node[anchor=north] at (axis cs:-0.4,0) {$y$};

    \node[stat221Red,anchor=west] at (axis cs:0.20,0.12) {tangent at $x$};
    \node[black,anchor=west] at (axis cs:0.72,0.72) {$f$};
  \end{axis}
\end{tikzpicture}


  \caption{Convexity first-order condition: a convex function lies above every tangent line/hyperplane.
  In one dimension, \eqref{eq:convex-first-order} reads $f(y)\ge f(x) + f'(x)(y-x)$.}
  \label{fig:convexity-tangent}
\end{figure}

\begin{definition}[$m$-strong convexity (via the Hessian)]
Let $f:S\to\R$ be twice differentiable.
We say $f$ is \emph{$m$-strongly convex} on $S$ if there exists $m>0$ such that
\begin{equation}
  \Hess f(x) \succeq m I
  \qquad \text{for all } x \in S.
  \label{eq:strong-convex-hess}
\end{equation}
\end{definition}

Here (and throughout), for symmetric matrices $A,B$ we write $A \succeq B$ (equivalently $B \preceq A$) to mean $A-B$ is positive semidefinite, i.e.\ $x^\top(A-B)x\ge 0$ for all $x\in\R^d$.

\begin{definition}[$M$-smoothness (Lipschitz gradient and Hessian bounds)]
Let $f:S\to\R$ be differentiable.
We say $f$ has \emph{$M$-Lipschitz gradient} (is \emph{$M$-smooth}) if
\begin{equation}
  \norm{\nabla f(x)-\nabla f(y)}_2 \le M \norm{x-y}_2
  \qquad \text{for all } x,y \in S.
  \label{eq:lipschitz-gradient}
\end{equation}
If $f$ is twice differentiable, then \eqref{eq:lipschitz-gradient} is equivalent to the operator-norm bound
\begin{equation}
  \norm{\Hess f(x)}_{\mathrm{op}} \le M
  \qquad \text{for all } x \in S.
  \label{eq:smooth-hess}
\end{equation}
For convex twice differentiable functions this is the same as $0 \preceq \Hess f(x) \preceq M I$.
More generally, the one-sided upper bound $\Hess f(x) \preceq M I$ is weaker than \eqref{eq:lipschitz-gradient}, but it is still enough for the descent lemma below.
\end{definition}

\begin{theorem}[Quadratic bounds from Hessian bounds]
\label{thm:quadratic-bounds}
Let $f:S\to\R$ be twice differentiable on a convex set $S$.
\begin{itemize}[leftmargin=*]
  \item If $\Hess f(z) \preceq MI$ for all $z\in S$, then for all $x,y\in S$,
  \[
    f(y) \le f(x) + \nabla f(x)^\top (y-x) + \frac{M}{2}\norm{y-x}_2^2.
  \]
  \item If $\Hess f(z) \succeq mI$ for all $z\in S$ (i.e., $f$ is $m$-strongly convex on $S$), then for all $x,y\in S$,
  \[
    f(y) \ge f(x) + \nabla f(x)^\top (y-x) + \frac{m}{2}\norm{y-x}_2^2.
  \]
\end{itemize}
In particular, if both bounds hold, then for all $x,y\in S$,
\begin{equation}
  f(x) + \nabla f(x)^\top (y-x) + \frac{m}{2}\norm{y-x}_2^2
  \le f(y)
  \le f(x) + \nabla f(x)^\top (y-x) + \frac{M}{2}\norm{y-x}_2^2.
  \label{eq:quad-bounds}
\end{equation}
\end{theorem}
\begin{proof}
Fix $x,y\in S$ and define $\phi(t)=f(x+t(y-x))$ for $t\in[0,1]$.
By the chain rule,
\[
\phi'(t)=\nabla f(x+t(y-x))^\top (y-x),
\qquad
\phi''(t)=(y-x)^\top \Hess f(x+t(y-x)) (y-x).
\]
If $\Hess f \preceq MI$ on $S$, then $\phi''(t) \le M\norm{y-x}_2^2$; if $\Hess f \succeq mI$ on $S$, then $\phi''(t) \ge m\norm{y-x}_2^2$.
Applying Taylor's theorem to $\phi(1)$ around $t=0$ with Lagrange remainder gives
\[
\phi(1)=\phi(0)+\phi'(0)+\tfrac12 \phi''(\xi)
\]
for some $\xi\in(0,1)$.
Plugging in the appropriate bound on $\phi''(\xi)$ yields the corresponding inequality; when both Hessian bounds hold we obtain \eqref{eq:quad-bounds}.
\end{proof}

The upper bound in \cref{thm:quadratic-bounds} is often called the \emph{descent lemma} (it is a quantitative version of ``$f$ is well-approximated by its tangent for small steps'').
The lower bound is the analytic content of strong convexity.
Together they are the basic calculus inequalities that let us prove convergence rates and justify step-size rules.

Geometrically, the inequalities in \eqref{eq:quad-bounds} say that, on $S$, the function $f$ is sandwiched between two quadratic bowls around any point $x$.
The lower bound (parameter $m$) prevents $f$ from being ``too flat'' and is what enables linear rates for first-order methods.
The upper bound (parameter $M$) prevents $f$ from having ``too much curvature'' and is what makes fixed (or backtracked) step sizes stable.

\begin{figure}[t]
  \centering
  \begin{tikzpicture}
  \begin{axis}[
    width=0.9\linewidth,
    height=5.6cm,
    xmin=-1.3, xmax=1.3,
    ymin=0.0, ymax=4.0,
    axis lines=left,
    ytick=\empty,
    xtick={0},
    xticklabels={$x^\star$},
    xticklabel style={font=\small,yshift=-0.2em},
    xtick style={stat221Gray!60},
    xlabel={$x$},
    xlabel style={at={(axis description cs:1,0)},anchor=north east,yshift=-0.35em},
    ylabel={$f(x)$},
    legend style={
      draw=stat221Gray!25,
      fill=white,
      fill opacity=0.90,
      text opacity=1,
      rounded corners,
      font=\footnotesize,
      at={(axis description cs:0.50,0.98)},
      anchor=north,
      inner sep=2.5pt,
    },
    legend cell align=left,
    clip=false,
  ]
    % A smooth strongly convex function with minimizer at x*=0.
    \addplot[black,thick,domain=-1.3:1.3,samples=220] {0.2 + x^2 + 0.2*x^4};
    \addlegendentry{$f$}

    % Quadratic upper/lower bounds at x* (illustrative).
    \addplot[stat221Red,thick,domain=-1.3:1.3,samples=220] {0.2 + 2*x^2}; % (M/2)||x-x*||^2 with M=4
    \addlegendentry{\textcolor{stat221Red}{$f(x^\star)+\tfrac{M}{2}(x-x^\star)^2$}}
    \addplot[stat221Blue,thick,domain=-1.3:1.3,samples=220] {0.2 + 0.5*x^2}; % (m/2)||x-x*||^2 with m=1
    \addlegendentry{\textcolor{stat221Blue}{$f(x^\star)+\tfrac{m}{2}(x-x^\star)^2$}}

    % Mark x*.
    \addplot[only marks,mark=*,mark size=1.6pt] coordinates {(0,0.2)};
    \addplot[densely dashed,stat221Gray!60] coordinates {(0,0) (0,0.2)};
  \end{axis}
\end{tikzpicture}

  \caption{If $mI \preceq \Hess f \preceq MI$ on a region containing $x^\star$, then $f$ is ``sandwiched'' between two quadratics around the minimizer:
  $f(x^\star)+\tfrac{m}{2}\norm{x-x^\star}^2 \le f(x)\le f(x^\star)+\tfrac{M}{2}\norm{x-x^\star}^2$.}
  \label{fig:quadratic-sandwich}
\end{figure}

\begin{corollary}[Uniqueness of the minimizer]
If $f$ is $m$-strongly convex on $S$, then $f$ has at most one minimizer on $S$.
\end{corollary}
\begin{proof}
Suppose $x^\star,y^\star\in S$ are two distinct minimizers, and define
\[
\phi(t)=f\bigl((1-t)x^\star + ty^\star\bigr), \qquad t\in[0,1].
\]
Then
\[
\phi''(t)=(y^\star-x^\star)^\top \Hess f\bigl((1-t)x^\star + ty^\star\bigr)(y^\star-x^\star)
\ge m\norm{y^\star-x^\star}_2^2 > 0,
\]
so $\phi$ is strictly convex on $[0,1]$.
Therefore
\[
f\Bigl(\frac{x^\star+y^\star}{2}\Bigr)=\phi\Bigl(\frac12\Bigr)
< \frac{\phi(0)+\phi(1)}{2}
= f(x^\star),
\]
which contradicts the minimality of $x^\star$ and $y^\star$.
\end{proof}

\begin{corollary}[Strong convexity links gradients, suboptimality, and distance]
\label{cor:sc-links}
Assume $f:\R^d\to\R$ is differentiable and $m$-strongly convex, and let $x^\star = \argmin_{x\in\R^d} f(x)$.
Then:
\begin{enumerate}[leftmargin=*,label=(\alph*)]
  \item $\nabla f(x^\star)=0$,
  \item for all $x \in \R^d$, $f(x) - f(x^\star) \le \frac{1}{2m}\norm{\nabla f(x)}_2^2$,
  \item for all $x \in \R^d$, $\norm{x-x^\star}_2 \le \frac{1}{m}\norm{\nabla f(x)}_2$.
\end{enumerate}
\end{corollary}
\begin{proof}
(a) is the first-order optimality condition for unconstrained differentiable convex minimization.

(b) Apply the left inequality in \cref{thm:quadratic-bounds} and minimize the resulting quadratic lower bound over $y\in\R^d$:
for any fixed $x$,
\[
f(y) \ge f(x) + \nabla f(x)^\top (y-x) + \frac{m}{2}\norm{y-x}_2^2.
\]
The right-hand side, viewed as a function of $y$, is minimized at $\hat y = x - \frac{1}{m}\nabla f(x)$.
Plugging in $\hat y$ gives
\[
f(y) \ge f(x) - \frac{1}{2m}\norm{\nabla f(x)}_2^2 \qquad \forall y,
\]
and choosing $y=x^\star$ yields (b).

(c) Strong convexity implies (see \cref{thm:quadratic-bounds} with $x=x^\star$ and $\nabla f(x^\star)=0$)
\[
f(x) \ge f(x^\star) + \frac{m}{2}\norm{x-x^\star}_2^2.
\]
Combine this with (b) to obtain $\frac{m}{2}\norm{x-x^\star}_2^2 \le \frac{1}{2m}\norm{\nabla f(x)}_2^2$, i.e. $\norm{x-x^\star}_2 \le \frac{1}{m}\norm{\nabla f(x)}_2$.
\end{proof}
For unconstrained smooth strongly convex objectives, these simple inequalities are extremely useful: they let us translate between three common notions of ``being close to optimum'' (small function gap, small gradient norm, and small distance to $x^\star$).
We will use them to turn a decrease in $f$ into a geometric (linear) convergence rate.
